% ============================================================
% MODULO 04 - OBJETIVOS
% ============================================================
\section{Objetivos}

\subsection{Pergunta norteadora}

O que a literatura nacional e internacional, de 2004 a 2026, descreve sobre a rela\c{c}\~ao entre brincar livre em espa\c{c}os externos ou na vizinhan\c{c}a e a atenua\c{c}\~ao da intensidade de sintomas de TDAH em pr\'e-escolares, e como esses achados dialogam com a cl\'inica humanista?

\subsection{Objetivo geral}

Sintetizar evid\^encias e lacunas sobre ambiente, brincar social e TDAH pr\'e-escolar, de modo a tornar intelig\'ivel a mudan\c{c}a comportamental observada no caso \'indice.

\subsection{Objetivos espec\'ificos}

\begin{enumerate}[label=\alph*),leftmargin=1.9cm]
  \item mapear as evid\^encias sobre natureza, espa\c{c}o verde e gravidade de sintomas de TDAH em crian\c{c}as;
  \item mapear as evid\^encias sobre brincar livre entre pares e autorregula\c{c}\~ao na faixa de 3 a 6 anos;
  \item articular os fundamentos da cl\'inica humanista -- Rogers e Axline -- que ajudam a explicar a tranquiliza\c{c}\~ao observada tamb\'em em sess\~ao.
\end{enumerate}
